% !TEX program = tectonic
\documentclass[preprint,12pt]{elsarticle}
\biboptions{numbers}

\usepackage{iftex}
\usepackage[T1]{fontenc}
\ifPDFTeX
  \usepackage[utf8]{inputenc}
\fi
\usepackage{lmodern}
\usepackage[english]{babel}
\usepackage{microtype}
\usepackage{amsmath,amssymb,amsthm}
\usepackage{booktabs,tabularx,multirow}
\usepackage{xcolor,graphicx,float,placeins}
\usepackage{tikz}
\usetikzlibrary{arrows.meta,positioning,fit,calc,shapes.geometric}
\usepackage{hyperref}
\hypersetup{colorlinks=true,linkcolor=black,citecolor=black,urlcolor=blue}
\urlstyle{same}
\usepackage[capitalise,noabbrev]{cleveref}
\usepackage{orcidlink}

\newcolumntype{Y}{>{\raggedright\arraybackslash}X}
\newcolumntype{L}[1]{>{\raggedright\arraybackslash}p{#1}}
\setlength{\emergencystretch}{3em}
\hfuzz=3pt
\setlength{\textfloatsep}{7pt plus 1pt minus 2pt}
\setlength{\floatsep}{6pt plus 1pt minus 1pt}
\setlength{\intextsep}{7pt plus 1pt minus 1pt}
\setcounter{topnumber}{5}
\setcounter{bottomnumber}{4}
\setcounter{totalnumber}{8}
\renewcommand{\topfraction}{0.95}
\renewcommand{\bottomfraction}{0.85}
\renewcommand{\textfraction}{0.05}
\renewcommand{\floatpagefraction}{0.75}
\raggedbottom

\newtheorem{definition}{Definition}[section]
\newtheorem{property}[definition]{Property}
\newtheorem{remark}[definition]{Remark}

\newcommand{\AFD}{\textsc{Aura File Defender}}
\newcommand{\Clean}{\mathsf{Clean}}
\newcommand{\Suspicious}{\mathsf{Suspicious}}
\newcommand{\Blocked}{\mathsf{Blocked}}
\newcommand{\Unprocessed}{\mathsf{Unprocessed}}
\newcommand{\Structural}{\mathsf{Structural}}
\newcommand{\Semantic}{\mathsf{Semantic}}
\newcommand{\Accept}{\mathsf{Accept}}
\newcommand{\Parse}{\mathsf{Parse}}
\newcommand{\Profile}{\mathsf{Profile}}
\newcommand{\Rebuild}{\mathsf{Rebuild}}
\newcommand{\Deliver}{\mathsf{Deliver}}

\IfFileExists{results/generated-results.tex}{%
  % Generated only from a schema-valid controlled server run.
% The pending values keep the drafting manuscript compilable without creating
% empirical claims before the release gate is satisfied.
\newcommand{\AFDRunID}{pending}
\newcommand{\AFDCommit}{pending}
\newcommand{\AFDCorpusSamples}{pending}
\newcommand{\AFDAcceptedSamples}{pending}
\newcommand{\AFDBlockedSamples}{pending}
\newcommand{\AFDFalseBlockingRate}{pending}
\newcommand{\AFDMedianLatency}{pending}
\newcommand{\AFDPninetyfiveLatency}{pending}
\newcommand{\AFDPeakMemory}{pending}
\newcommand{\AFDFuzzBudget}{pending}
\newcommand{\AFDFidelitySummary}{pending}
%
}{%
  \newcommand{\AFDRunID}{pending}%
  \newcommand{\AFDCommit}{pending}%
  \newcommand{\AFDCorpusSamples}{pending}%
  \newcommand{\AFDAcceptedSamples}{pending}%
  \newcommand{\AFDBlockedSamples}{pending}%
  \newcommand{\AFDFalseBlockingRate}{pending}%
  \newcommand{\AFDMedianLatency}{pending}%
  \newcommand{\AFDPninetyfiveLatency}{pending}%
  \newcommand{\AFDPeakMemory}{pending}%
  \newcommand{\AFDFuzzBudget}{pending}%
  \newcommand{\AFDFidelitySummary}{pending}%
}

\date{}

\begin{document}
\begin{frontmatter}

\title{Aura File Defender: Endpoint-Local Canonical Reconstruction of Media Attachments in End-to-End Encrypted Messaging}

\author{Oleksandr Melnychenko\texorpdfstring{\,\orcidlink{0000-0001-8565-7092}}{}\texorpdfstring{\corref{cor1}}{}}
\ead{melnychenko@khmnu.edu.ua}
\cortext[cor1]{Corresponding author.}
\address{Khmelnytskyi National University, Khmelnytskyi, Ukraine}
\journal{Journal of Information Security and Applications}

\begin{abstract}
End-to-end encryption protects attachment confidentiality in transit but does
not determine whether decrypted media should be opened at an endpoint. This
paper presents \AFD, an endpoint-local reconstruction and validation method for
media attachments. The method places the same decision boundary before
outgoing encryption and after incoming decryption, so the relay neither sees
plaintext nor decides acceptance. A file is deliverable only when its portable
name, caller-supplied media type, byte-derived format profile, required
reconstruction level, and independently parsed canonical output converge.
Structural and semantic reconstruction are distinguished explicitly, and a
policy that requires semantic processing cannot fall back to copied codec
payloads. Complete output consumption prevents trailing-object acceptance, while
non-media, ambiguous, malformed, over-budget, and non-clean results expose no
deliverable bytes. The evaluation protocol separates carrier-level malware
properties from behavioral malware classification and preregisters corpus,
differential-parser, fault-injection, fuzzing, fidelity, and performance gates.
Empirical values are withheld from this drafting version until a pinned server
run satisfies the complete machine-readable release gate.
\end{abstract}

\begin{keyword}
content disarm and reconstruction \sep end-to-end encryption \sep canonical parsing \sep media security \sep polyglot files \sep malware carriers
\end{keyword}

\end{frontmatter}

\section{Introduction}\label{sec:introduction}

Encrypted messaging changes where attachment inspection can occur. A relay
that is not permitted to decrypt an attachment cannot reconstruct its content,
whereas an endpoint that opens the decrypted object inherits the parser and
decoder attack surface. Transport confidentiality and file admissibility are
therefore separate security decisions. The first protects bytes in transit.
The second determines which representation, if any, may reach a renderer. This
distinction complements established secure-messaging models, which analyze
communication confidentiality and key evolution but do not certify the safety
of a decrypted attachment representation~\cite{unger2015,cohngordon2020}.

Conventional scanning and content reconstruction are often described as
gateway services. That placement does not preserve the confidentiality model
of a protected conversation. It also encourages a binary interpretation in
which a file is either labelled malicious or passed through. Media
reconstruction supports a narrower and testable statement: an accepted output
belongs to a constrained grammar and was derived at the reconstruction strength
required by policy. This statement does not imply that retained pixels or audio
samples are semantically harmless.

The research question is whether endpoint-local reconstruction can preserve a
blind relay while establishing a reproducible, format-specific acceptance
boundary for attachments. The contribution is a composed method rather than a
new codec. It combines concrete input-profile convergence, an explicit
reconstruction-strength order, independent complete output parsing, fail-closed
publication, and a preregistered evidence protocol. These elements are evaluated
as separate claims so that a successful parser test is not presented as malware
detection and a successful decode is not presented as canonical validation.

\section{Related Work and Claim Boundary}\label{sec:related}

File-processing discrepancies are an established evasion surface. Jana and
Shmatikov showed that disagreements between detector-side and host-side format
interpretation can permit malware to evade analysis without changing its
executable behavior~\cite{jana2012}. Language-theoretic security consequently
treats accepted input as a language that should be specified explicitly and
recognized before use~\cite{momot2016}. Nail operationalizes a related idea by
deriving parsers and generators from one grammar~\cite{bangert2014}. EverParse
adds machine-checked safety, inverse parsing and serialization, and unique
representations for authenticated message formats~\cite{ramananandro2019}.
Interval parsing grammars address random-access and type--length--value
dependencies that arise in binary file formats~\cite{zhang2023}. These systems
motivate grammar-centred validation, but they do not define an attachment
publication policy for protected messaging endpoints.

Polyglot research demonstrates why a leading signature or one successful
decoder is insufficient. Albertini et al. systematize format combinations in
the different setting of authenticated encryption without key
commitment~\cite{albertini2022}. Koch et al. document polyglot attack chains,
evaluate format identification, and present an image sanitization
tool~\cite{koch2025}. You et al. show that even implementations nominally
parsing the same ZIP format can disagree at the semantic level~\cite{you2025}.
These findings support differential validation and complete-consumption checks.
They do not justify treating any single detector as an oracle.

The contribution here is neither a new media codec nor a universal malware
classifier. It is an enforcement composition whose deployment point and
acceptance rule are jointly specified. The method places reconstruction at both
ends of a blind encrypted relay. It requires name, declared media type, and byte
evidence to converge on a concrete profile before dispatch. It then requires
the performed reconstruction strength and a separate complete parser to agree
before any bytes can be published. The evaluation binds each of these claims to
a distinct machine-readable gate rather than inferring security from aggregate
scanner accuracy.

The behavioral terms used in \cref{tab:carrier-map} follow the conventional
malware categories summarized by NIST~\cite{nist80083}. The measured properties
remain narrower. They concern carrier grammar, embedded active objects,
metadata, and reconstruction outcome rather than program behavior after
execution.

\begin{table}[t]
\centering
\caption{Permitted interpretation of malware terminology in the evaluation.}
\label{tab:carrier-map}
\begin{tabularx}{\textwidth}{L{0.22\textwidth}YY}
\toprule
Behavioral label & Measured attachment property & Permitted conclusion \\
\midrule
File-infector virus & Executable prefix, modified host grammar, or appended executable body & Carrier rejected by type, grammar, or polyglot policy \\
Worm or downloader & Script, executable, shortcut, package, or active document & Unsupported active carrier rejected \\
Trojan & Executable or script presented with media evidence & Classifier disagreement or unsupported grammar rejected \\
Ransomware or spyware & Executable carrier or specified privacy metadata & Carrier rejected or named metadata removed \\
Exploit-bearing media & Reproducible malformed structure or decoder regression & Frozen path rejected or reconstructed this sample \\
Steganographic payload & Information retained in pixels or samples & Outside the reconstruction claim \\
\bottomrule
\end{tabularx}
\end{table}

\section{System and Adversary Model}\label{sec:model}

The adversary controls attachment bytes, display name, declared media type,
container metadata, dimensions, counts, offsets, codec payload, and request
frequency. The policy, authenticated application binary, and local publication
directory are trusted. External decoders are trusted only within the stated
process-isolation assumptions. Compromise of the endpoint before processing and
harmful semantic content visible in a legitimate raster or waveform are outside
the attachment-reconstruction claim.

\begin{figure}[t]
\centering
\resizebox{\textwidth}{!}{%
\begin{tikzpicture}[
  node distance=5mm and 7mm,
  box/.style={draw,rounded corners=2pt,align=center,minimum height=8mm,minimum width=22mm,font=\small},
  guard/.style={box,fill=blue!7},
  crypto/.style={box,fill=green!8},
  relay/.style={box,fill=gray!12},
  arrow/.style={-{Latex[length=2mm]},thick}
]
\node[box] (source) {source\\attachment};
\node[guard,right=of source] (outguard) {classify\\rebuild\\validate};
\node[crypto,right=of outguard] (encrypt) {encrypt};
\node[relay,right=of encrypt] (relay) {blind relay\\ciphertext only};
\node[crypto,right=of relay] (decrypt) {decrypt};
\node[guard,right=of decrypt] (inguard) {classify\\rebuild\\validate};
\node[box,right=of inguard] (render) {renderer\\or block};
\draw[arrow] (source) -- (outguard);
\draw[arrow] (outguard) -- node[above,font=\scriptsize]{clean only} (encrypt);
\draw[arrow] (encrypt) -- (relay);
\draw[arrow] (relay) -- (decrypt);
\draw[arrow] (decrypt) -- (inguard);
\draw[arrow] (inguard) -- node[above,font=\scriptsize]{clean only} (render);
\end{tikzpicture}
}
\caption{Symmetric endpoint placement of the reconstruction boundary. The
relay transports ciphertext and does not authorize attachment delivery.}
\label{fig:placement}
\end{figure}

\section{Canonical Reconstruction Method}\label{sec:method}

Let $x$ denote input bytes, $n$ a display name, $m$ a caller-supplied media
type, and $p$ a policy. Three independent functions derive concrete format
profiles from the final extension, parsed media type, and bounded byte evidence.
The media type is not treated as ground truth. HTTP semantics likewise warn
that media-type sniffing can create security risks and that byte inspection
cannot determine every processing model~\cite{rfc9110}:
\begin{equation}
  e(x,n,m)=\bigl(E(n),M(m),B(x)\bigr).
  \label{eq:evidence}
\end{equation}
Aliases are represented explicitly. Compatibility does not collapse distinct
grammars merely because they belong to the same top-level media family.

\begin{definition}[Reconstruction strength]
The performed strength belongs to
$\{\Unprocessed,\Structural,\Semantic\}$ with
\begin{equation}
  \Unprocessed \prec \Structural \prec \Semantic.
  \label{eq:strength}
\end{equation}
Structural reconstruction emits a newly parsed container but may retain codec
payload. Semantic reconstruction decodes media state and encodes a new object.
\end{definition}

For candidate output $y=\Rebuild_{p,k}(x)$, an independent parser returns a
profile, component count, and terminal offset. Delivery is authorized by the
conjunction
\begin{align}
\Accept(p,x,n,m,y,k)={}&N_p(n) \land C_p(e,k) \land L_p(k) \notag\\
 &\land \Parse_k(y) \land T_k(y)=|y| \notag\\
 &\land K_p(k,y) \land S_p(y) \land Q_p(x,y),
\label{eq:acceptance}
\end{align}
where $N_p$ is portable-name acceptance, $C_p$ is profile convergence, $L_p$
is reconstruction-level satisfaction, $K_p$ is canonical component closure,
$S_p$ is the post-reconstruction signature policy, and $Q_p$ contains resource
bounds. No single classifier or decoder exit status can satisfy
\cref{eq:acceptance}.

\begin{property}[Fail-closed publication]
For every result $r$,
\begin{equation}
  \Deliver(r)=1 \quad\Longleftrightarrow\quad r.\mathit{verdict}=\Clean.
  \label{eq:delivery}
\end{equation}
Both $\Suspicious$ and $\Blocked$ carry zero deliverable bytes. Filesystem
publication is separated from reconstruction and occurs only through a
same-directory temporary file after a complete write.
\end{property}

\begin{figure}[t]
\centering
\resizebox{0.96\textwidth}{!}{%
\begin{tikzpicture}[
  node distance=6mm and 6mm,
  box/.style={draw,rounded corners=2pt,align=center,minimum height=9mm,minimum width=25mm,font=\small},
  check/.style={box,fill=blue!7},
  transform/.style={box,fill=green!8},
  terminal/.style={box,fill=gray!12},
  arrow/.style={-{Latex[length=2mm]},thick}
]
\node[check] (evidence) {profile\\convergence};
\node[transform,right=of evidence] (rebuild) {required-strength\\reconstruction};
\node[check,right=of rebuild] (parse) {independent\\complete parse};
\node[check,right=of parse] (policy) {resource and\\signature policy};
\node[terminal,right=of policy] (publish) {clean-only\\publication};
\draw[arrow] (evidence) -- (rebuild);
\draw[arrow] (rebuild) -- (parse);
\draw[arrow] (parse) -- (policy);
\draw[arrow] (policy) -- (publish);
\end{tikzpicture}
}
\caption{Conjunctive acceptance path. Failure or uncertainty at any stage
produces no deliverable bytes. No stage can authorize publication alone.}
\label{fig:acceptance-path}
\end{figure}

\section{Format Profiles and Polyglot Handling}\label{sec:formats}

The output language is an allowlist of registered canonical grammars. Each
parser starts at offset zero, rejects budget exhaustion as failure, and reports
complete consumption. Prefix camouflage, appended grammar, overlapping
grammar, and nested active objects are measured separately. Semantic
neutralization is claimed only when no source container or compressed payload
bytes are copied into the output and the new representation passes an
independent complete parser. Normative format rules supply the lower-level
closure conditions. For example, the PNG specification fixes signature, chunk,
and terminal-IEND requirements~\cite{png3}, while the Ogg Opus specification
defines mandatory headers and page organization~\cite{rfc7845}.

\begin{table}[t]
\centering
\caption{Release treatment by media family. Exact accepted profiles are frozen
in the machine-readable registry.}
\label{tab:profiles}
\begin{tabularx}{\textwidth}{L{0.20\textwidth}L{0.25\textwidth}Y}
\toprule
Family & Principal transformation & Acceptance boundary \\
\midrule
PNG, JPEG, static WebP & Pixel decode and re-encode & Complete canonical raster parse and decode \\
GIF & Frame decode and re-encode & Canvas, frame, timing, trailer, and complete sub-block parse \\
APNG & Bounded chunk reconstruction & Structural profiles only. Semantic-required profiles block \\
Animated WebP & Bounded RIFF and frame reconstruction & Metadata-free structural profile, complete frame parse and decode. Semantic-required profiles block \\
MP3, WAV, FLAC & Metadata-aware structural reconstruction or transcode & Complete frame/chunk grammar and integrity fields \\
Ogg/Opus & Page and packet reconstruction or transcode & CRC, sequence, packet, codec, and EOS closure \\
MP4/M4A & Parent-aware atom reconstruction or transcode & Media-type alias convergence, Apple advisory-atom removal, track and sample-table closure \\
WebM & Nested EBML reconstruction or transcode & DocType, track, block, lacing, timestamp, and boundary closure \\
Other formats & No pass-through & Non-media and unsupported classes block \\
\bottomrule
\end{tabularx}
\end{table}

\section{Implementation}\label{sec:implementation}

The reference implementation is a Rust workspace with separate policy,
signature, observability, foreign-function, media-handler, and facade crates.
Input-profile convergence precedes handler dispatch. Every handler reports the
performed reconstruction strength and exact structural removal counters where
the parser can identify source regions. Output acceptance uses a separate
validator module rather than trusting a handler declaration.

Animated WebP is routed separately from static WebP. Its structural handler
removes ICC, EXIF, and XMP chunks, bounds canvas and frame work, and retains
constrained compressed frame payloads. The independent validator reparses the
complete RIFF object and decodes every frame, but the performed level remains
structural because the frame payloads are not regenerated. For ISO BMFF audio,
\texttt{audio/mp4}, \texttt{audio/x-m4a}, and \texttt{audio/m4a} converge on
one profile. The rewriter removes excluded Apple advisory atoms before the
same complete sample-table and media-extent validation used for other MP4
outputs.

External FFmpeg and FFprobe work is supervised with a wall-clock deadline,
cancellation token, bounded diagnostics, output-file monitoring, and process
termination on all exit paths. Linux execution adds a dedicated process group,
no-new-privileges state, disabled core dumps, CPU and output-file limits, and a
bounded address space. The C and Swift boundaries repeat the clean-only output
invariant. The Apple artifact contains iOS, Simulator, Mac Catalyst, and macOS
slices built from one locked source revision. Its build remaps local source
paths, fixes archive epochs, and canonically serializes the XCFramework property
list before distribution comparison.

\section{Evaluation Protocol}\label{sec:evaluation}

The corpus is stratified into canonical benign media, benign noncanonical
media, malformed structures, classifier disagreement, polyglot carriers,
resource-exhaustion cases, unsupported active classes, and decoder regressions.
Ground truth requires a reproducible generator, two independent reference
parsers, or documented adjudication by at least two reviewers. Every accepted
output is saved by digest and checked by independent parser families.

For deterministic structural or lossless paths, stability requires
\begin{equation}
  \Rebuild_{p,k}(\Rebuild_{p,k}(x))=\Rebuild_{p,k}(x).
  \label{eq:idempotence}
\end{equation}
Lossy semantic paths use the weaker named property of canonical closure and are
evaluated with separately frozen fidelity thresholds. Performance observations
are grouped by format and reconstruction level. Confidence intervals use a
hierarchical bootstrap over server sessions and corpus samples, while corpus
order is deterministically permuted for each session.

\begin{table}[t]
\centering
\caption{Controlled-run summary. Values remain pending until the run manifest
reports a passed scientific release gate.}
\label{tab:results}
\begin{tabular}{ll}
\toprule
Run identifier & \AFDRunID \\
Source commit & \AFDCommit \\
Corpus samples & \AFDCorpusSamples \\
Clean / non-clean outcomes & \AFDAcceptedSamples{} / \AFDBlockedSamples \\
False-blocking rate & \AFDFalseBlockingRate \\
Median / p95 latency & \AFDMedianLatency{} / \AFDPninetyfiveLatency \\
Peak resident memory & \AFDPeakMemory \\
Fuzzing budget & \AFDFuzzBudget \\
Fidelity summary & \AFDFidelitySummary \\
\bottomrule
\end{tabular}
\end{table}

\section{Discussion and Limitations}\label{sec:discussion}

Canonical reconstruction narrows the representation delivered to a renderer.
It does not prove the absence of malicious meaning or information encoded in
retained semantic content. Structural reconstruction cannot claim removal of a
payload inside copied codec frames. Semantic reconstruction depends on the
decoder, encoder, operating-system isolation, and the fidelity thresholds used
in the evaluation. Coverage-guided fuzzing and parser agreement are empirical
evidence rather than a proof of parser correctness.

The method also makes an explicit availability trade-off. Unknown formats,
missing evidence, and unavailable required tools are non-deliverable under the
release profile. The false-blocking analysis must therefore be reported by
concrete format and corpus stratum rather than hidden in one aggregate rate.
Future work may add a semantic animated-WebP path, stronger operating-system
sandboxing for external decoders, and typed export of all parser resource
counters.

\section{Conclusion}\label{sec:conclusion}

\AFD{} treats attachment admissibility as an endpoint-local, format-specific
decision that remains compatible with a blind encrypted relay. Its central
claim is not universal malware detection. It is the enforceable conjunction of
profile convergence, required reconstruction strength, complete canonical
output parsing, bounded execution, and clean-only publication. The frozen
evaluation protocol is designed to show precisely which part of that
conjunction is supported by each experiment.

\section*{Data and Code Availability}

\enlargethispage{4\baselineskip}

The source revision, policy, classifier registry, corpus manifest, raw
observations, analysis scripts, tool digests, and run-manifest hashes will be
reported with the evaluated artifact. Raw malicious samples that cannot be
redistributed will remain in the isolated corpus and will be represented by
provenance, mutation recipes, and cryptographic digests.

\begin{thebibliography}{99}

\bibitem{unger2015}
N. Unger, S. Dechand, J. Bonneau, S. Fahl, H. Perl, I. Goldberg, and
M. Smith, ``SoK: Secure Messaging,'' in \emph{2015 IEEE Symposium on
Security and Privacy}, 2015, pp. 232--249.
\url{https://doi.org/10.1109/SP.2015.22}

\bibitem{cohngordon2020}
K. Cohn-Gordon, C. Cremers, B. Dowling, L. Garratt, and D. Stebila,
``A formal security analysis of the Signal messaging protocol,''
\emph{Journal of Cryptology}, vol. 33, pp. 1914--1983, 2020.
\url{https://doi.org/10.1007/s00145-020-09360-1}

\bibitem{jana2012}
S. Jana and V. Shmatikov, ``Abusing file processing in malware detectors
for fun and profit,'' in \emph{2012 IEEE Symposium on Security and Privacy},
2012, pp. 80--94. \url{https://doi.org/10.1109/SP.2012.15}

\bibitem{momot2016}
F. Momot, S. Bratus, S. M. Hallberg, and M. L. Patterson, ``The seven
turrets of Babel: A taxonomy of LangSec errors and how to expunge them,''
in \emph{2016 IEEE Cybersecurity Development}, 2016, pp. 45--52.
\url{https://doi.org/10.1109/SecDev.2016.019}

\bibitem{bangert2014}
J. Bangert and N. Zeldovich, ``Nail: A practical tool for parsing and
generating data formats,'' in \emph{11th USENIX Symposium on Operating
Systems Design and Implementation}, 2014, pp. 615--628.
\url{https://www.usenix.org/conference/osdi14/technical-sessions/presentation/bangert}

\bibitem{ramananandro2019}
T. Ramananandro, A. Delignat-Lavaud, C. Fournet, N. Swamy, T. Chajed,
N. Kobeissi, and J. Protzenko, ``EverParse: Verified secure zero-copy
parsers for authenticated message formats,'' in \emph{28th USENIX Security
Symposium}, 2019, pp. 1465--1482.
\url{https://www.usenix.org/conference/usenixsecurity19/presentation/delignat-lavaud}

\bibitem{zhang2023}
J. Zhang, G. Morrisett, and G. Tan, ``Interval parsing grammars for file
format parsing,'' \emph{Proceedings of the ACM on Programming Languages},
vol. 7, no. PLDI, art. 150, 2023.
\url{https://doi.org/10.1145/3591264}

\bibitem{albertini2022}
A. Albertini, T. Duong, S. Gueron, S. K{\"o}lbl, A. Luykx, and S. Schmieg,
``How to abuse and fix authenticated encryption without key commitment,''
in \emph{31st USENIX Security Symposium}, 2022, pp. 3291--3308.
\url{https://www.usenix.org/conference/usenixsecurity22/presentation/albertini}

\bibitem{koch2025}
L. Koch, S. Oesch, A. Sadovnik, B. Weber, A. Chaulagain, M. Dixson,
J. Dixon, M. Huettel, C. Watson, J. Hartman, and R. Patulski, ``On the
abuse and detection of polyglot files,'' in \emph{Proceedings of the ACM Web
Conference 2025}, 2025, pp. 4810--4822.
\url{https://doi.org/10.1145/3696410.3714814}

\bibitem{you2025}
Y. You, J. Chen, Q. Wang, and H. Duan, ``My ZIP isn't your ZIP: Identifying
and exploiting semantic gaps between ZIP parsers,'' in \emph{34th USENIX
Security Symposium}, 2025, pp. 431--450.
\url{https://www.usenix.org/conference/usenixsecurity25/presentation/you}

\bibitem{nist80083}
M. Souppaya and K. Scarfone, \emph{Guide to Malware Incident Prevention
and Handling for Desktops and Laptops}, NIST Special Publication 800-83
Revision 1, 2013. \url{https://doi.org/10.6028/NIST.SP.800-83r1}

\bibitem{rfc9110}
R. Fielding, M. Nottingham, and J. Reschke, \emph{HTTP Semantics},
RFC 9110, 2022. \url{https://www.rfc-editor.org/rfc/rfc9110}

\bibitem{png3}
W3C PNG Working Group, \emph{Portable Network Graphics (PNG)
Specification, Third Edition}, W3C Recommendation, 2025.
\url{https://www.w3.org/TR/png-3/}

\bibitem{rfc7845}
T. Terriberry, R. Lee, and R. Giles, \emph{Ogg Encapsulation for the Opus
Audio Codec}, RFC 7845, 2016.
\url{https://www.rfc-editor.org/rfc/rfc7845}

\end{thebibliography}

\end{document}
